\begin{lemma}
\label{lemma-computation}
In Situation \ref{situation-affine}. Assume $D$ and $\Omega_D$ are as in
(\ref{equation-D}) and (\ref{equation-omega-D}).
Let $\lambda \in D$. Let $D'$ be the $p$-adic completion of
$$
D[z]\langle \xi \rangle/(\xi - (z^p - \lambda))
$$
and let $\Omega_{D'}$ be the $p$-adic completion of the module of
divided power differentials of $D'$ over $A$. For any pair $(M, \nabla)$
over $D$ satisfying (\ref{item-complete}), (\ref{item-connection}),
(\ref{item-integrable}), and (\ref{item-topologically-quasi-nilpotent})
the canonical map of complexes (\ref{equation-base-change-map-complexes})
$$
c_M^\bullet : M \otimes_D^\wedge \Omega^\bullet_D
\longrightarrow
M \otimes_D^\wedge \Omega^\bullet_{D'}
$$
has the following property: There exists a map $e_M^\bullet$
in the opposite direction such that both $c_M^\bullet \circ e_M^\bullet$
and $e_M^\bullet \circ c_M^\bullet$ are homotopic to multiplication by $p$.
\end{lemma}

\begin{proof}
We will prove this using Lemma \ref{lemma-find-homotopy} with $a = p$.
Thus we have to find $\theta : D' \to D'$ and prove
(\ref{item-d-a-zero}), (\ref{item-direct-sum}), (\ref{item-theta-linear}),
(\ref{item-integrate}), (\ref{item-injective}), (\ref{item-factor}),
(\ref{item-horizontal}). We first collect some information about the rings
$D$ and $D'$ and the modules $\Omega_D$ and $\Omega_{D'}$.

\medskip\noindent
Writing
$$
D[z]\langle \xi \rangle/(\xi - (z^p - \lambda))
=
D\langle \xi \rangle[z]/(z^p - \xi - \lambda)
$$
we see that $D'$ is the $p$-adic completion of the free $D$-module
$$
\bigoplus\nolimits_{i = 0, \ldots, p - 1}
\bigoplus\nolimits_{n \geq 0}
z^i \xi^{[n]} D
$$
where $\xi^{[0]} = 1$.
It follows that $D \to D'$ has a continuous $D$-linear section, in particular
$D \to D'$ is universally injective, i.e., (\ref{item-injective}) holds.
We think of $D'$ as a divided power algebra
over $A$ with divided power ideal $\overline{J}' = \overline{J}D' + (\xi)$.
Then $D'$ is also the $p$-adic completion of the divided power envelope
of the ideal generated by $z^p - \lambda$ in $D$, see
Lemma \ref{lemma-describe-divided-power-envelope}. Hence
$$
\Omega_{D'} = \Omega_D \otimes_D^\wedge D' \oplus D'\text{d}z
$$
by Lemma \ref{lemma-module-differentials-divided-power-envelope}.
This proves (\ref{item-direct-sum}). Note that (\ref{item-d-a-zero})
is obvious.

\medskip\noindent
At this point we construct $\theta$. (We wrote a PARI/gp script theta.gp
verifying some of the formulas in this proof which can be found in the
scripts subdirectory of the Stacks project.) Before we do so we compute
the derivative of the elements $z^i \xi^{[n]}$. We have
$\text{d}z^i = i z^{i - 1} \text{d}z$. For $n \geq 1$ we have
$$
\text{d}\xi^{[n]} =
\xi^{[n - 1]} \text{d}\xi =
- \xi^{[n - 1]}\text{d}\lambda + p z^{p - 1} \xi^{[n - 1]}\text{d}z
$$
because $\xi = z^p - \lambda$. For $0 < i < p$ and $n \geq 1$ we have
\begin{align*}
\text{d}(z^i\xi^{[n]})
& =
iz^{i - 1}\xi^{[n]}\text{d}z + z^i\xi^{[n - 1]}\text{d}\xi \\
& =
iz^{i - 1}\xi^{[n]}\text{d}z + z^i\xi^{[n - 1]}\text{d}(z^p - \lambda) \\
& =
- z^i\xi^{[n - 1]}\text{d}\lambda +
(iz^{i - 1}\xi^{[n]} + pz^{i + p - 1}\xi^{[n - 1]})\text{d}z \\
& =
- z^i\xi^{[n - 1]}\text{d}\lambda +
(iz^{i - 1}\xi^{[n]} + pz^{i - 1}(\xi + \lambda)\xi^{[n - 1]})\text{d}z \\
& =
- z^i\xi^{[n - 1]}\text{d}\lambda +
((i + pn)z^{i - 1}\xi^{[n]} + p\lambda z^{i - 1}\xi^{[n - 1]})\text{d}z
\end{align*}
the last equality because $\xi \xi^{[n - 1]} = n\xi^{[n]}$.
Thus we see that
\begin{align*}
\partial_z(z^i) & = i z^{i - 1} \\
\partial_z(\xi^{[n]}) & = p z^{p - 1} \xi^{[n - 1]} \\
\partial_z(z^i\xi^{[n]}) & =
(i + pn) z^{i - 1} \xi^{[n]} + p \lambda z^{i - 1}\xi^{[n - 1]}
\end{align*}
Motivated by these formulas we define $\theta$ by the rules
$$
\begin{matrix}
\theta(z^j)
& = & p\frac{z^{j + 1}}{j + 1}
& j = 0, \ldots p - 1, \\
\theta(z^{p - 1}\xi^{[m]})
& = & \xi^{[m + 1]}
& m \geq 1, \\
\theta(z^j \xi^{[m]})
& = &
\frac{p z^{j + 1} \xi^{[m]} - \theta(p\lambda z^j \xi^{[m - 1]})}{(j + 1 + pm)}
& 0 \leq j < p - 1, m \geq 1
\end{matrix}
$$
where in the last line we use induction on $m$ to define our choice of
$\theta$. Working this out we get (for $0 \leq j < p - 1$ and $1 \leq m$)
$$
\theta(z^j \xi^{[m]}) =
\textstyle{\frac{p z^{j + 1} \xi^{[m]}}{(j + 1 + pm)} -
\frac{p^2 \lambda z^{j + 1} \xi^{[m - 1]}}{(j + 1 + pm)(j + 1 + p(m - 1))} +
\ldots +
\frac{(-1)^m p^{m + 1} \lambda^m z^{j + 1}}
{(j + 1 + pm) \ldots (j + 1)}}
$$
although we will not use this expression below. It is clear that $\theta$
extends uniquely to a $p$-adically continuous $D$-linear map on $D'$.
By construction we have (\ref{item-theta-linear}) and (\ref{item-integrate}).
It remains to prove (\ref{item-factor}) and (\ref{item-horizontal}).

\medskip\noindent
Proof of (\ref{item-factor}) and (\ref{item-horizontal}).
As $\theta$ is $D$-linear and continuous it suffices to prove that
$p - \theta \circ \partial_z$,
resp.\ $(\theta \otimes 1) \circ \text{d}_1 - \text{d}_1 \circ \theta$
gives an element of $D$, resp.\ $\Omega_D$ when evaluated on the
elements $z^i\xi^{[n]}$\footnote{This can be done by direct computation:
It turns out that $p - \theta \circ \partial_z$ evaluated on
$z^i\xi^{[n]}$ gives zero except for $1$ which is mapped to $p$ and
$\xi$ which is mapped to $-p\lambda$. It turns out that 
$(\theta \otimes 1) \circ \text{d}_1 - \text{d}_1 \circ \theta$
evaluated on $z^i\xi^{[n]}$ gives zero except for $z^{p - 1}\xi$
which is mapped to $-\lambda$.}.
Set $D_0 = \mathbf{Z}_{(p)}[\lambda]$ and
$D_0' = \mathbf{Z}_{(p)}[z, \lambda]\langle \xi \rangle/(\xi - z^p + \lambda)$.
Observe that each of the expressions above is an element of
$D_0'$ or $\Omega_{D_0'}$. Hence it suffices to prove the result
in the case of $D_0 \to D_0'$. Note that $D_0$ and $D_0'$
are torsion free rings and that $D_0 \otimes \mathbf{Q} = \mathbf{Q}[\lambda]$
and $D'_0 \otimes \mathbf{Q} = \mathbf{Q}[z, \lambda]$.
Hence $D_0 \subset D'_0$ is the subring of elements annihilated
by $\partial_z$ and (\ref{item-factor})
follows from (\ref{item-integrate}), see the discussion directly preceding
Lemma \ref{lemma-find-homotopy}. Similarly, we have
$\text{d}_1(f) = \partial_\lambda(f)\text{d}\lambda$ hence
$$
\left((\theta \otimes 1) \circ \text{d}_1 - \text{d}_1 \circ \theta\right)(f)
=
\left(\theta(\partial_\lambda(f)) - \partial_\lambda(\theta(f))\right)
\text{d}\lambda
$$
Applying $\partial_z$ to the coefficient we obtain
\begin{align*}
\partial_z\left(
\theta(\partial_\lambda(f)) - \partial_\lambda(\theta(f))
\right)
& =
p \partial_\lambda(f) - \partial_z(\partial_\lambda(\theta(f))) \\
& =
p \partial_\lambda(f) - \partial_\lambda(\partial_z(\theta(f))) \\
& =
p \partial_\lambda(f) - \partial_\lambda(p f) = 0
\end{align*}
whence the coefficient does not depend on $z$ as desired.
This finishes the proof of the lemma.
\end{proof}